\documentclass{article}

\usepackage[utf8]{inputenc}
\usepackage[german]{babel}
\usepackage{listings}
\usepackage{xcolor}

\lstset{
    basicstyle=\ttfamily\small,
    breaklines=true,
    backgroundcolor=\color{gray!10}
}

\begin{document}
    \begin{center}
        \Large \textbf{Huntsman Challenge 2: Konfigurationsmanagement in großem Maßstab} \\ [6pt]
        \normalsize
        Autor: Maximilian Jakob Maag B.Sc. \\
        Version: 1.0
    \end{center}

    \section{Intro}
    Diese Challenge befasst sich mit der zentralen Verwaltung von Konfigurationen auf einer großen Anzahl von Servern.
    Ihr Ziel ist es, eine konsistente und wartbare Konfiguration aller verwalteten Server sicherzustellen.
    Sie werden lernen, wie man Konfigurationsänderungen skaliert durchführt und dabei die Sicherheit und Stabilität gewährleistet.

    \section{Aufgaben}

    \subsection{Aufgabe 1}
    Der Rabbit hat 20+ Nextcloud Server in verschiedenen Regionen deployt.
    Sie müssen diese Server in Ihrem Inventory organisieren und über Gruppen verwalten können.
    \\ \\
    Erstellen Sie Hostgruppen nach Kriterien wie Region, Umgebung und Service-Typ.

    \subsection{Aufgabe 2}
    Erstellen Sie ein Basis-Konfigurationsplaybook, das auf alle neu hinzugekommenen Server angewandt wird.
    Dieses Playbook sollte folgende Aufgaben durchführen:
    \begin{itemize}
        \item System Packages aktualisieren
        \item Zeitzone setzen
        \item NTP konfigurieren
        \item Basis-Sicherheitsrichtlinien anwenden (Firewall, SSH Konfiguration)
        \item Monitoring Agent installieren
        \item Logging konfigurieren
    \end{itemize}

    \subsection{Aufgabe 3}
    Erstellen Sie regionspeziefische Playbooks für die 20+ Nextcloud Server.
    \\ \\
    Nutzen Sie das Playbook um sicherzustellen, dass:
    \begin{itemize}
        \item Alle Server die gleiche PHP Version haben
        \item PostgreSQL auf der gleichen Version läuft
        \item Apache/Nginx Konfigurationen konsistent sind
        \item SSL Zertifikate aktuell sind
        \item Alle Server auf dem gleichen Patchlevel sind
    \end{itemize}

    \subsection{Aufgabe 4}
    Implementieren Sie ein Verifizierungs-Playbook, das nach jeder Konfigurationsänderung durchgeführt wird.
    Das Playbook soll überprüfen:
    \begin{itemize}
        \item Connectivity zu allen Hosts
        \item Verfügbare Ressourcen (Disk Space, Memory, CPU)
        \item Laufende Services und deren Status
        \item Korrektheit der Konfigurationsdateien
        \item Sicherheitsbaselines (offene Ports, User Accounts, SSH Keys)
    \end{itemize}

    \subsection{Aufgabe 5}
    Erstellen Sie ein Playbook zur koordinierten Anwendung von Sicherheits-Patches.
    \\ \\
    Besonderheiten:
    \begin{itemize}
        \item Patches sollten Gruppen-weise durchgeführt werden (z.B. zuerst Staging, dann Production)
        \item Es sollte Health Checks vor und nach jedem Update geben
        \item Automatische Rollback Mechanismen sollten vorhanden sein
        \item Services sollten mit minimaler Ausfallzeit neu gestartet werden
    \end{itemize}

    \subsection{Aufgabe 6}
    Erstellen Sie ein Performance Monitoring Playbook, das regelmäßig folgende Metriken erfasst:
    \begin{itemize}
        \item CPU, Memory, Disk Auslastung
        \item Netzwerk Traffic
        \item Service Response Times
        \item Fehlerquoten in Log Files
    \end{itemize}
    
    Die erfassten Daten sollen in eine zentrale Monitoring Lösung exportiert werden (z.B. Prometheus, ELK Stack).

    \subsection{Aufgabe 7}
    Dokumentieren Sie alle Ihre Playbooks und erstellen Sie einen Runbook mit häufigen Aufgaben.
    Der Runbook sollte enthalten:
    \begin{itemize}
        \item Wie man neue Server zum Inventory hinzufügt
        \item Wie man ein Config Update durchführt
        \item Wie man Probleme diagnostiziert
        \item Wie man Server aus der Verwaltung entfernt
    \end{itemize}

\end{document}
